\documentclass[12pt]{article}
\newcommand{\bottomMargin}{2cm}

\newcommand{\header}{
	ДИЗАЙН И АНАЛИЗ НА АЛГОРИТМИ - \\ ПРАКТИКУМ\\
	\vspace{0.1cm}
	Летен семестър, 2026 г., второ контролно\\
	\vspace{0.1cm}
}
\newcommand{\tl}{$0.4$ сек.}
\newcommand{\ml}{$256$ MB}

\input{structure.tex}

\begin{document}

\section{Задача K7. Пътища}

В страната $X$ има $N$ на брой града. Между някои двойки градове съществуват директни пътища, като мрежата от пътища е свързана - т.е. от всеки град може да се стигне до всеки друг (евентуално минавайки през междинни градове).

В рамките на мащабна инфраструктурна реформа покрай Giro-то правителството взима следното решение: всички съществуващи директни пътища се премахват, а на тяхно място се изграждат директни пътища между \emph{всички двойки} градове, между които \emph{не е имало} директен път при старата мрежа. Един от съветниците забелязва, че след тази промяна е възможно да съществува двойка градове, между които да не съществува път(макар и не пряк). Това накарало господа
министрите сериозно да се замислят и те възложили на група от експерти да пресметнат при изграждане на новата шосейна мрежа какъв минимален брой от старите преки шосета да се запазят и ремонтират, така че да няма двойка
населени места, между които да липсва път.

Напишете програма \textbf{roads}, която да помогне на експертите да се справят с поставената им задача.

\subsection{Вход}

На първия ред на стандартния вход се намират две естествени числа $N$ и $M$,
разделени с интервал:
\begin{itemize}
	\item $N$ - брой на градовете;
	\item $M$ - брой на старите директни пътища.
\end{itemize}

Следват $M$ реда, на всеки от които са дадени две естествени числа $a$ и $b$ ($a \neq b$), разделени с интервал, означаващи, че между градове $a$ и $b$ съществува стар директен път. Градовете са номерирани от $1$ до $N$.

\subsection{Изход}

На стандартния изход отпечатайте едно цяло число - минималния брой стари директни пътища, които трябва да бъдат запазени.

\subsection{Ограничения}

\begin{itemize}
	\item $2 \leq N \leq 1000$
	\item $1 \leq M \leq 500\,000$
	\item $1 \leq a, b \leq N$, $a \neq b$
\end{itemize}

\subsection{Примери}

\begin{table}[ht]
	\begin{tblr}{|X[30,l]|X[10,l]|X[60,l]|}
		\hline
		\textbf{Вход} & \textbf{Изход} & \textbf{Обяснение} \\
		\hline
		\texttt{4 4\\1 2\\2 3\\3 4\\4 1} &
		\texttt{1} &
		{Старата мрежа е цикъл $1$–$2$–$3$–$4$–$1$. Новите пътища, които ще бъдат добавени, са $1$–$3$ и $2$–$4$. Достатъчно е да се запази \textbf{едно} старо ребро, например $1$–$2$, за да стане мрежата свързана.} \\
		\hline
		\texttt{3 3\\1 2\\2 3\\1 3} &
		\texttt{2} &
		{Между всички двойки градове има стари пътища. Достатъчно е да запазим $2$ от тях, например $2-3$ и $3-1$, за да има път от всеки град до всеки друг.} \\
		\hline
	\end{tblr}
\end{table}

\end{document}
